% vim: tw=78 ai spell spelllang=en

\documentclass[10pt]{llncs}
\usepackage{color}
\usepackage{url}
\pagestyle{plain}

\newcommand\Color[2]{\textcolor{#1}{#2}}
\newcommand{\SM}{{\textit{SM}}}

\begin{document}

\title{Checking Properties Described by State Machines: On Synergy of
Instrumentation, Slicing, and Symbolic Execution}
\author{Jiri~Slaby, Jan~Strej\v{c}ek, and Marek~Trt\'{\i}k}
\institute{Faculty of Informatics, Masaryk University\\
  Botanick\'a 68a, 60200 Brno, Czech Republic\\
  \email{\{slaby,strejcek,trtik\}@fi.muni.cz}}

\maketitle

\begin{abstract}
  We introduce a novel technique for checking properties described by finite
  state machines. The technique is based on a synergy of three well-known
  methods: instrumentation, program slicing, and symbolic execution. More
  precisely, we generate a code that corresponds to state machine transitions.
  This code is then instrumented into a given program. The error states are
  represented as \texttt{assert} calls. At this point we know the slicing
  criteria. They are exactly the conditions of the assertions. Hence we can
  slice the program to reduce its size without affecting runs of state
  machines. The slicing takes 
  
  And then we
  symbolically execute the sliced program to find real violations of
  the checked properties, i.e.~real bugs. Depending on the kind of symbolic
  execution, the technique can be applied as a stand-alone bug finding
  technique, or to weed out some false positives from an output of another
  bug-finding tool. We provide several examples demonstrating the practical
  applicability of our technique.
\end{abstract}

\end{document}
